\documentclass[10pt,twocolumn]{article}
\usepackage[T1]{fontenc}
\usepackage[utf8]{inputenc}
\usepackage{lmodern}
\usepackage[margin=1.8cm,top=2.4cm,bottom=2cm,columnsep=0.6cm]{geometry}
\usepackage{fancyhdr}
\usepackage{titlesec}
\usepackage{booktabs}
\usepackage{amsmath,amssymb,amsfonts}
\usepackage{microtype}
\usepackage{xcolor}
\usepackage{mdframed}
\usepackage{parskip}
\usepackage{enumitem}
\usepackage{hyperref}

\definecolor{rulered}{RGB}{180,30,30}
\definecolor{boxgray}{RGB}{245,245,245}
\definecolor{keywordgray}{RGB}{100,100,100}

\pagestyle{fancy}
\fancyhf{}
\fancyhead[L]{\textsc{\small Claude Mag}}
\fancyhead[C]{\small\textit{Machine Learning Research Digest}}
\fancyhead[R]{\small 2026-08-03}
\fancyfoot[C]{\small\thepage}
\renewcommand{\headrulewidth}{0.8pt}

\titleformat{\section}{\normalsize\bfseries\color{rulered}}{}{0pt}{}%
  [\vspace{-4pt}\color{rulered}\rule{\columnwidth}{0.4pt}]
\titlespacing{\section}{0pt}{8pt}{4pt}

\hypersetup{colorlinks=true,urlcolor=rulered,linkcolor=black}

\begin{document}
\setlength{\parindent}{0pt}
\setlength{\parskip}{4pt}

{\small\color{keywordgray}\textbf{Keywords:}
  code generation \textbullet{}
  reinforcement learning \textbullet{}
  runtime efficiency \textbullet{}
  execution feedback}\par\vspace{4pt}

{\LARGE\bfseries\raggedright
  RLPF: Reinforcement Learning from Performance
  Feedback for Code Generation}\par\vspace{4pt}

{\large\itshape\raggedright
  A Staged Reward Trains Code Models to Prefer Faster Correct
  Programs---Raising Correct-and-Efficient Rate from 8.1\% to 38.6\%}\par\vspace{6pt}

{\small\textbf{By} Huihao Jing, Haozhe Cui, Wenbin Hu, Shaojin Chen,
  Haochen Shi, Changxuan Fan, Yuxuan Liu, Hanyu Yang, Sirui Zhang,
  Ziyi Chen, Haoran Li, Yangqiu Song
  (Hong Kong Univ.\ of Science and Technology)}\par\vspace{2pt}

{\small\color{keywordgray}arXiv:2607.27271 \textbullet{} July 29, 2026}
\par\vspace{2pt}\noindent{\color{rulered}\rule{\columnwidth}{0.6pt}}\vspace{6pt}

Code models trained with execution feedback reward correctness---not
efficiency. Two programs can pass the same test suite while differing by
orders of magnitude in runtime. \textbf{RLPF} (Reinforcement Learning
from Performance Feedback) introduces a staged reward: correctness is
gated first; passing programs then receive a speedup signal relative to a
reference solution. On a competitive programming benchmark, RLPF raises a
base model from 8.1\% to 38.6\% CGRE (Correct-and-Graded Relative
Efficiency), closing two-thirds of the gap to a frontier model.

\begin{mdframed}[backgroundcolor=boxgray,linecolor=rulered,linewidth=1pt,
    innerleftmargin=6pt,innerrightmargin=6pt,
    innertopmargin=6pt,innerbottommargin=6pt]
\textbf{\small AT A GLANCE}\par\vspace{4pt}
\begin{description}[leftmargin=0pt,labelsep=4pt,itemsep=2pt,parsep=0pt,
    font=\small\bfseries]
  \item[Problem:] Correctness-only RL for code leaves no signal for runtime
    efficiency; programs that pass tests may be asymptotically slow.
  \item[Why it matters:] Systems code, competitive programming, and
    latency-sensitive applications require efficient implementations as a
    primary objective.
  \item[Method:] Two-stage reward: correctness gate first; relative
    speedup over a reference solution for programs that pass.
  \item[Result:] Correct-and-runnable rate 11.1\%$\to$54.6\%; CGRE
    8.1\%$\to$38.6\%.
  \item[Evidence:] Staged reward outperforms correctness-only RL by
    7.3 CGRE points; validated across array, sorting, graph, and DP tasks.
\end{description}
\end{mdframed}

\section{The Big Picture}

The code generation field has converged on a compelling training paradigm:
train a language model on execution feedback. Sample programs from the
model, execute them against a test suite, and use the pass/fail signal as
a reward for reinforcement learning. This has driven dramatic improvements
on correctness-centric benchmarks such as HumanEval, MBPP, and LiveCodeBench.

But correctness is not the only property that matters. In systems
programming, competitive programming, and any latency-sensitive application,
a correct but slow program is often not good enough. Runtime efficiency
should be a first-class training objective. RLPF is the first study that
makes efficiency a primary RL training signal for code generation, rather
than an evaluation-time metric.

\section{Why Existing Approaches Fall Short}

Correctness-only RL has a critical gap: the model learns to produce any
program that passes the test suite, regardless of runtime. For most
benchmarks, the test suite uses small inputs where even $O(n^3)$
solutions pass within the time limit. The model thus learns to write
correct but asymptotically slow code.

Three properties of runtime make it a difficult RL reward:

\textbf{1. Fragility.} Runtime is only meaningful after a program compiles
and passes all correctness checks. If 80\% of sampled programs fail to
compile, a pure efficiency reward provides near-zero gradient.

\textbf{2. Non-comparability across tasks.} An absolute runtime of
100\,ms may be excellent for one problem and catastrophic for another.
A relative metric is required.

\textbf{3. Mode collapse risk.} Optimizing for efficiency without a
correctness constraint causes the model to produce fast-but-wrong programs
(e.g., always returning a hardcoded answer).

RLPF addresses all three with a staged reward design.

\section{Core Mechanism}

RLPF uses PPO with a two-stage reward function:

\textbf{Stage 1 --- Correctness gate.} A sampled program $\hat{c}$ is
executed against the full test suite. If it fails to compile, fails any
test, or times out, it receives zero reward and the performance stage is
skipped.

\textbf{Stage 2 --- Performance reward.} If $\hat{c}$ passes all tests,
its runtime $t(\hat{c})$ is measured on a distribution of inputs. The
reward is the relative speedup over a reference correct solution:
\[
  r_{\text{perf}}(\hat{c}) = \frac{t_{\text{ref}} - t(\hat{c})}{t_{\text{ref}}}
\]
This is bounded in $(-\infty, 1]$, with 0 meaning same speed as the
reference and 1 meaning infinite speedup.

\section{Technical Formulation}

The combined reward is:
\[
  r(\hat{c}) = \alpha \cdot r_{\text{correct}} +
  (1 - \alpha) \cdot r_{\text{perf}}(\hat{c}) \cdot
  \mathbf{1}\bigl[r_{\text{correct}} = 1\bigr]
\]
where $\alpha$ is annealed from 1.0 to 0.3 over training via a cosine
schedule.

The PPO objective maximizes:
\[
  \mathcal{J}(\theta) = \mathbb{E}\!\left[
    \min\!\left(\rho_t A_t,\;
      \mathrm{clip}(\rho_t, 1{-}\epsilon, 1{+}\epsilon) A_t
    \right)
  \right] - \beta D_{\mathrm{KL}}(\pi_\theta \| \pi_{\mathrm{ref}})
\]
where $\rho_t = \pi_\theta(\hat{c}|p)/\pi_{\mathrm{old}}(\hat{c}|p)$
and $A_t$ is the advantage from $r(\hat{c})$.

The evaluation metric CGRE is:
\[
  \text{CGRE} = \frac{1}{|P|}\sum_{p\in P}
  \mathbf{1}[\hat{c}_p\text{ passes}] \cdot
  \max\!\!\left(0,\; 1 - \frac{t(\hat{c}_p)}{t_{\mathrm{ref}}(p)}\right)
\]
This rewards both correctness and efficiency: failing programs
contribute 0, and faster programs contribute proportionally more.

\section{Learning or Inference Procedure}

Training uses a dataset of competitive programming problems covering
arrays, sorting, graph algorithms, and dynamic programming. Reference
solutions are taken from accepted competitive submissions.

PPO is run for 3,000 steps with batch size 32, sampling 8 rollouts per
problem per step. Each rollout is compiled and executed; the staged reward
is computed online. The $\alpha$ annealing follows a cosine curve,
ensuring the correctness stage dominates early training when the model
produces mostly incorrect programs.

The KL divergence penalty $D_{\mathrm{KL}}(\pi_\theta \| \pi_{\mathrm{ref}})$
relative to the initial model prevents mode collapse onto a few fast-but-narrow
templates.

\section{What the Guarantee Says}

The $\alpha$ annealing provides a correct-first guarantee: during the first
phase of training ($\alpha$ near 1), the policy is incentivized exclusively
to produce correct programs. By the time $\alpha$ decreases and the
performance reward becomes significant, the correctness rate is already high
enough that most sampled programs enter the performance stage, providing
dense efficiency gradients.

\section{Experimental Findings}

\begin{center}\small
\begin{tabular}{lcc}
\toprule
Method & Correct-and-Run & CGRE \\
\midrule
Base model & 11.1\% & 8.1\% \\
Correctness-only RL & 47.3\% & 18.4\% \\
RLPF (staged) & \textbf{54.6\%} & \textbf{38.6\%} \\
GPT-5.4 (zero-shot) & 58.2\% & 41.1\% \\
\bottomrule
\end{tabular}
\end{center}

RLPF raises CGRE by 30.5 points over the base model and by 20.2 points
over correctness-only RL. The gap to a frontier closed-source model is
reduced to 2.5 CGRE points.

\textbf{By problem category:} Largest CGRE gains occur on array and DP
problems (44.2\% CGRE), where efficient algorithmic patterns (prefix sums,
memoization) are learnable. Smaller gains on graph problems (28.1\%),
where efficiency requires sophisticated data structures harder to induce
from supervision.

\section{Ablations and Interpretation}

\textbf{Reward annealing:} Fixing $\alpha = 0.5$ throughout reduces CGRE
by 8.2 points, as the efficiency signal disrupts early correctness learning
and produces more incorrect programs.

\textbf{Reference quality:} Using a 10th-percentile runtime reference
(suboptimal) reduces CGRE by 4.1 points---reference quality is a primary
bottleneck and should be set to the best known solution.

\textbf{KL penalty:} Removing $D_{\mathrm{KL}}$ causes collapse onto 3--5
high-efficiency code templates, reducing CGRE diversity and failing on
problems that do not fit those templates. Performance drops by 6.3 CGRE.

\textbf{Number of rollouts:} 8 rollouts per problem per step is the
sweet spot; 4 rollouts degrades CGRE by 3.1 points due to lower variance
in the reward signal.

\section*{Reference}

Huihao Jing, Haozhe Cui, Wenbin Hu, Shaojin Chen, Haochen Shi,
Changxuan Fan, Yuxuan Liu, Hanyu Yang, Sirui Zhang, Ziyi Chen,
Haoran Li, Yangqiu Song.
\textbf{RLPF: Reinforcement Learning from Performance Feedback for Code
Generation}.
arXiv:2607.27271, July 29, 2026.
\url{https://arxiv.org/abs/2607.27271}

\end{document}
